\chapter{Kapitel 1}

Miss Brooke had that kind of beauty which seems to be thrown into relief by poor dress. Her hand and wrist were so finely formed that she could wear sleeves not less bare of style than those in which the Blessed Virgin appeared to Italian painters; and her profile as well as her stature and bearing seemed to gain the more dignity from her plain garments, which by the side of provincial fashion gave her the impressiveness of a fine quotation from the Bible,--or from one of our elder poets,--in a paragraph of to-day's newspaper. She was usually spoken of as being remarkably clever, but with the addition that her sister Celia had more common-sense. Nevertheless, Celia wore scarcely more trimmings; and it was only to close observers that her dress differed from her sister's, and had a shade of coquetry in its arrangements; for Miss Brooke's plain dressing was due to mixed conditions, in most of which her sister shared. The pride of being ladies had something to do with it: the Brooke connections, though not exactly aristocratic, were unquestionably "good:" if you inquired backward for a generation or two, you would not find any yard-measuring or parcel-tying forefathers--anything lower than an admiral or a clergyman; and there was even an ancestor discernible as a Puritan gentleman who served under Cromwell, but afterwards conformed, and managed to come out of all political troubles as the proprietor of a respectable family

\pex
\begingl
\gla Konta Brok lodar en slåg sla brikkhæd, känn mykken enmer tilvådför par slyste håg.//
\glb Miss Brooke have.\Pst{} the type of beauty, what make.\Pass{} more noticeable by:the simple clothing//
\glft Miss Brooke had that kind of beauty which seems to be thrown into relief by poor dress.//
\endgl
\xe
\pex
\begingl
